\section{Introduction}
\label{sec:introduction}

\subsection{Motivation}
The effective use of lighting is paramount in creating immersive environments for stage performances, film production, and artistic installations.
The rapid advancement of \ac{LED} technology and embedded systems has democratized access to high-quality lighting tools.
However, a gap often remains between rigid, high-cost professional equipment and fragile, feature-limited hobbyist solutions.
The RLCV3 (\ac{RGB} \ac{LED} Controller Version 3) project was initiated to address this specific niche.
It stems from the need for a highly customizable, robust, and modular lighting ecosystem capable of integrating seamlessly into professional workflows via standard protocols, while remaining accessible and adaptable for custom "light sculpture" configurations.
Building upon the lessons learned from previous iterations, specifically the limitations in mechanical modularity and control resolution encountered in Version 2, this project aims to deliver a professional-grade solution suitable for \ac{DIY} fabrication.

\subsection{Problem Statement and Objectives}
The primary challenge addressed by this work is the development of a unified control architecture capable of driving distinct lighting form factors—specifically linear tubes, flat panels, and spotlights—without requiring disparate hardware designs.
Existing solutions often lack the flexibility to switch between wired \ac{DMX} and wireless Art-Net control within a compact form factor.
Furthermore, mechanical integration in custom builds is frequently overlooked, leading to thermal issues or structural instability.

The specific objectives of this project are:
\begin{itemize}
	\item To design a versatile embedded controller based on the ESP32-C3 \ac{MCU} that supports both wired DMX512 and wireless Art-Net communication.
	\item To develop firmware that adapts to multiple hardware profiles (Addressable \ac{RGB} strips, \ac{PWM} \ac{MOSFET} control, and \ac{I2C}-driven panels) through a centralized configuration mechanism.
	\item To engineer a robust mechanical enclosure system utilizing \ac{FDM} 3D printing and aluminum profiles that facilitates modular interconnection (the "Screw it!" system) and ensures effective thermal management.
	\item To implement a responsive local user interface using an \ac{OLED} display and rotary encoder for standalone configuration.
	\item To validate the system's performance in terms of thermal stability, power reliability, and software responsiveness suitable for live environments.
\end{itemize}

\subsection{Structure of the Article}
The remainder of this document is organized as follows:
\begin{itemize}
	\item \textit{Section \ref{sec:theoretical_background} \nameref{sec:theoretical_background}} provides the theoretical underpinnings, covering embedded systems, \ac{LED} technology, and communication protocols.
	\item \textit{Section \ref{sec:system_architecture} \nameref{sec:system_architecture}} presents the high-level hardware and software design, including the block diagrams and menu structures.
	\item \textit{Section \ref{sec:hardware_development} \nameref{sec:hardware_development}} details the electronic circuit design, \ac{PCB} fabrication, and component selection.
	\item \textit{Section \ref{sec:mechanical_design} \nameref{sec:mechanical_design}} describes the mechanical engineering aspects, focusing on the enclosure design, 3D printing materials, and the modular connection system.
	\item \textit{Section \ref{sec:software_development} \nameref{sec:software_development}} discusses the firmware implementation, including the event-driven architecture, library management, and user interface logic.
	\item \textit{Section \ref{sec:device_assebly} \nameref{sec:device_assebly}} outlines the practical steps for assembling the devices and flashing the firmware.
	\item Finally, \textit{Section \ref{sec:results} \nameref{sec:results}} analyzes the performance of the developed hardware and software, evaluating thermal behavior and overall system functionality.
\end{itemize}